\begin{abstract}
We check whether tools (search, an answer sheet or code execution) help a small language model reject a user's false claim when the model already knows the correct answer without pressure. Sycophancy means giving up a correct answer to agree with the user. This is an ethical risk when the model is used as a medical or education assistant, and people may think that a small model with tools is safer than a model that only uses its own knowledge. We test three models (Qwen2.5-3B, Llama-3.2-3B and Phi-4-mini, all 4-bit) on 146 medical-exam (MedQA) and 215 mathematics (MATH) questions that all three first answer correctly. In the same conversation the user then suggests a wrong answer and repeats it with stronger pressure, so there are two pressure levels. Three conditions are common to both datasets: no tool (T0), BM25 answer sheet with two entries (T1) and the exact answer-sheet entry (Exact). We also add dataset-specific tools: web search and textbook passages for MedQA, and a Python execution tool for MATH. Without a tool, Qwen keeps the correct MATH answer in only \mathQwenPtwoTzero\% of expert turns. Python execution raises expert-turn resistance by \deltaQwenPyPtwo{} to \deltaLlamaPyPtwo{} points for all three models, and the exact answer keeps Phi-4 at \medqaPhiPtwoExact\% on MedQA, but web snippets and textbook passages lower Phi-4's resistance from \medqaPhiPtwoTzero\% to \medqaPhiPtwoWeb\% and \medqaPhiPtwoText\%. Qwen rarely calls the tool because it has a position bias in the tool menu, and Llama gives the user's answer even when the tool gives it the correct answer. So using a tool does not always mean less sycophancy; tools must give direct evidence, and models must learn to weigh it against the user.
\end{abstract}
